\section{Conclusion}
\label{sec:conclusion}

Agents interleave, and the classical answer to interleaving assumes plans that
agents do not have. We identified the resulting problem, admitting a plan
revision after part of the plan has committed, and characterized it
completely: a revision can be made safe by regrouping the future if and only
if sealing the future certifies, and when it does not, neither a finer
regrouping nor waiting helps. We built a runtime that enforces the hypotheses
of the certificate instead of assuming them, and showed that the locks inside
a piece can be ordered by the same priority that repair-based protocols use
for notifications, which removes deadlock without giving up short holds.

With real agents, certified chopping is as safe as whole-workflow locking and
safer than repair on the workloads where repair depends on the model; with
priority-ordered locking it is also as fast as repair. Its failures leave
state clean rather than corrupted. And it has limits that we measured rather
than assumed away: when a read-modify-write straddles the reasoning gap the
right answer is not to chop, and when an agent escalates into shared state
after committing, the right answer is usually to refuse.

\paragraph{Artifact.}
The agent runtime, the eight protocols, the twelve cells, the brute-force
verification suite, every trial record of both model runs, and the scripts
that produce every table and figure in \S\ref{sec:eval:setup}--\S\ref{sec:eval:revisions}
are available in the accompanying repository. The durable runtime and its
verifier reproduce Table~\ref{tab:verify} and \S\ref{sec:eval:mechanism}
using only the Python standard library.
